\section{Introduction}

Autonomous agentic AI systems, pipelines of LLM-based agents that plan, execute, document, verify, and sometimes deploy their own work, create a practical version of a classical epistemic problem. When one agent produces code and another agent reviews it, neither judgment supplies an unquestionable ground of correctness. The verification regress, known in philosophy as Agrippa's trilemma or M\"unchhausen's trilemma \cite{sextus_empiricus,albert1985}, therefore becomes an engineering question: how should an agentic system be structured when no component can be treated as an infallible verifier?

Most current work on LLM agents studies performance under a fixed architecture: a given role decomposition, prompting strategy, tool interface, or benchmark task is evaluated for short-horizon success \cite{li2023camel,wu2024autogen,hong2023metagpt,qian2024chatdev,yang2024sweagent}. Much less is known about \emph{architectural dynamics}: how role structure changes during prolonged autonomous operation, which pressures trigger that change, and whether the resulting dynamics can be modeled rather than narrated post hoc. This paper treats architecture itself as the dependent variable.

We investigate three research questions. \textbf{RQ1}: What recurrent patterns of architectural change appear when an autonomous agentic system operates for multiple days without human intervention? \textbf{RQ2}: When multiple instances begin from nearly null task specification, which aspects of their developmental trajectory converge and which diverge? \textbf{RQ3}: Can the autonomy boundary between agentic systems and human overseers be formalized in a way that supports measurement, learning, and later empirical testing?

To answer these questions, we combine two empirical studies with a formal and conceptual analysis. The \emph{Wallfacer} study is a seven-day longitudinal observation of an autonomous software-engineering pipeline that passed through four architectural configurations. The \emph{Ralph} study is a three-run replication experiment in which independent instances started from empty repositories with no seed task. From these studies we extract two empirical regularities. First, we observe \emph{bottleneck migration}: improvements in one system dimension repeatedly relocate the binding constraint to another. Second, we observe convergence-then-divergence: independent runs converge on the same task family while diverging in implementation language, specialization, and feature strategy.

The paper makes four contributions, ordered by epistemic strength. \textbf{Empirical}: we document bottleneck migration in Wallfacer and convergence-then-divergence in Ralph. \textbf{Analytic}: we propose a \emph{role-fission} hypothesis stating that architectural differentiation becomes likely when conflicting cognitive modes overrun a shared context budget. \textbf{Formal}: we model trust calibration as Gaussian Process preference learning with uncertainty-aware escalation and frame architecture search as constrained preferential Bayesian optimization \cite{chu2005,gonzalez2017}. \textbf{Programmatic}: we connect the framework to the literatures on self-adaptive software, organizational design, cybernetics, and security, and we mark missing evidence with explicit TODOs rather than treating speculative claims as established results.

Several scope conditions are important. The empirical base is small: one longitudinal Wallfacer trajectory and three Ralph runs, all on a single LLM substrate and within software-engineering tasks. The architectural transitions in Wallfacer were adaptive interventions, not randomized treatments, so our causal language is intentionally cautious. Accordingly, the paper should be read as a hypothesis-generating systems study plus a formal research agenda, not as a definitive theory of all agentic systems.

The remainder of the paper proceeds from observation to conjecture. Sections~2 and~3 report the Wallfacer and Ralph experiments. Section~4 extracts the role-fission principle and its falsifiable implications. Section~5 formalizes bottleneck migration, trust calibration, and architecture search. Sections~6 and~7 extend the framework into more speculative territory while keeping epistemic status explicit. Sections~8 and~9 position the work relative to prior literature, enumerate threats to validity and missing experiments, and conclude.
